\section{Introduction}

For a positive integer $n$, let $j(n)$ be the least positive integer $m$
such that every interval of $m$ consecutive integers contains an integer
coprime to $n$. Write $\omega(n)$ for the number of distinct prime divisors
of $n$, with $\omega(1)=0$, and put
\[
 h(k)=\sup_{\substack{n\ge1\\\omega(n)\le k}}j(n),\qquad k\ge1.
\]
The starting point of the interval is arbitrary, and replacing $n$ by
the product of its distinct prime divisors does not change $j(n)$.
Thus $h(k)-1$ is the greatest length of an interval that can be covered
by the divisibility classes of at most $k$ primes.

Jacobsthal studied this function in a series of papers beginning in
1960 \cite{Jacobsthal1960}. Erd\H{o}s subsequently recorded Jacobsthal's question
whether $h(k)\ll k^2$ \cite[p.~163, equations (1) and (3)]{Erdos1962}.
In Erd\H{o}s's notation, $C(k)+1$ is the maximum of $j(n)$ over integers
with exactly $k$ distinct prime divisors. This equals our $h(k)$:
adjoining prime divisors cannot decrease $j(n)$.
We prove the following bound, which gives an affirmative answer and a
further iterated-logarithmic saving.

\begin{theorem}\label{thm:main}
There is an absolute constant $C>0$ such that
\[
 h(k)\le C\frac{k^2}{(\log\log(3k))^2}
 \qquad\text{for every integer }k\ge1.
\]
\end{theorem}

All logarithms are natural. In particular the theorem gives a
uniform quadratic bound for arbitrary sets of prime divisors, with
no restriction on their sizes or on the interval's translation.

